% GENERATED FILE. Edit canonical lessons, then run npm run generate:books.
% canonical-source-hash: fnv1a64:aca442dd172e41f6
% canonical-lessons: ML-C399-samrajyam, ML-C399-kotta, ML-C399-sainyam, ML-C399-pattalakkaran, ML-C399-kilatakkuka

\chapter{The Empire and Its Army}
\label{ch:ml-the-empire-and-its-army}
\hlchaptermodality{399}

\begin{chapteropening}
\textbf{By the end of this chapter:} \emph{I can say an empire, a fort, an army, a soldier and to conquer.}

Complete the last lesson of chapter 399: I can say an empire, a fort, an army, a soldier and to conquer.
\end{chapteropening}

\section[\texorpdfstring{\ml{സാമ്രാജ്യം} (sāmrājyaṁ)}{സാമ്രാജ്യം (sāmrājyaṁ)}]{\ml{സാമ്രാജ്യം} (sāmrājyaṁ) — an empire}

\label{lesson:ML-C399-samrajyam}



\begin{quote}
Before the new one: say the Malayalam for a sovereign, then the Malayalam for to reign.
\end{quote}

\subsection*{You'll want to know: \ml{സാമ്രാജ്യം}}
\textbf{\ml{സാമ്രാജ്യം}} — \emph{sāmrājyaṁ} — \textquotedblleft{}an empire\textquotedblright{}.

\begin{quote}
\ml{ആ} \ml{സാമ്രാജ്യം} \ml{ഒരുപാട്} \ml{രാജ്യങ്ങളിൽ} \ml{പരന്നു} \ml{കിടന്നു}.
\end{quote}

\begin{quote}
\emph{ā sāmrājyaṁ orupāṭŭ rājyaṅṅaḷil parannu kiṭannu.}
\end{quote}

\begin{quote}
That empire stretched over many countries.
\end{quote}

One ruler over many kingdoms.

\begin{cousinweb}
\textbf{\ml{സാമ്രാജ്യം}} (\emph{sāmrājyaṁ}) is borrowed from Sanskrit \emph{sāmrājya}, \textquotedblleft{}complete or universal sovereignty, empire\textquotedblright{} (Monier-Williams).
\end{cousinweb}

\begin{grammarlens}[title={What you've built}]
One more word for telling what happened.
\end{grammarlens}

\subsection*{Your turn}
\begin{itemize}
\raggedright
  \item \emph{Say it:} \emph{sāmrājyaṁ}
  \item \emph{Say it:} \emph{sāmrājyaṁ}, once more
  \item \emph{Say it:} \emph{ā sāmrājyaṁ orupāṭŭ rājyaṅṅaḷil parannu kiṭannu.} — That empire stretched over many countries.
  \item \emph{Recall:} say the Malayalam for a sovereign, then the Malayalam for to reign, then say \emph{sāmrājyaṁ} again
\end{itemize}

\begin{tcolorbox}[breakable,colback=teal!4,colframe=teal!35!black,title={Before you move on}]
What does \ml{സാമ്രാജ്യം} mean? (\textquotedblleft{}An empire\textquotedblright{}.) Say it once more.
\end{tcolorbox}
\section[\texorpdfstring{\ml{കോട്ട} (kōṭṭa)}{കോട്ട (kōṭṭa)}]{\ml{കോട്ട} (kōṭṭa) — a fort}

\label{lesson:ML-C399-kotta}



\begin{quote}
Before the new one: say the Malayalam for to reign, then the Malayalam for an empire.
\end{quote}

\subsection*{You'll want to know: \ml{കോട്ട}}
\textbf{\ml{കോട്ട}} — \emph{kōṭṭa} — \textquotedblleft{}a fort\textquotedblright{}.

\begin{quote}
\ml{കടൽത്തീരത്ത്} \ml{ഒരു} \ml{പഴയ} \ml{കോട്ട} \ml{ഉണ്ട്}.
\end{quote}

\begin{quote}
\emph{kaṭalttīrattŭ oru paḻaya kōṭṭa uṇṭŭ.}
\end{quote}

\begin{quote}
There is an old fort on the seashore.
\end{quote}

Plural \ml{കോട്ടകൾ}.

\begin{cousinweb}
\textbf{\ml{കോട്ട}} (\emph{kōṭṭa}) is native Dravidian: DEDR 2207 lists Malayalam \emph{kōṭṭa}, \textquotedblleft{}fort, residence\textquotedblright{}, beside Tamil \emph{kōṭṭai}, \textquotedblleft{}fort, castle\textquotedblright{}.
\end{cousinweb}

\begin{grammarlens}[title={What you've built}]
One more word for telling what happened.
\end{grammarlens}

\subsection*{Your turn}
\begin{itemize}
\raggedright
  \item \emph{Say it:} \emph{kōṭṭa}
  \item \emph{Say it:} \emph{kōṭṭa}, once more
  \item \emph{Say it:} \emph{kaṭalttīrattŭ oru paḻaya kōṭṭa uṇṭŭ.} — There is an old fort on the seashore.
  \item \emph{Recall:} say the Malayalam for to reign, then the Malayalam for an empire, then say \emph{kōṭṭa} again
\end{itemize}

\begin{tcolorbox}[breakable,colback=teal!4,colframe=teal!35!black,title={Before you move on}]
What does \ml{കോട്ട} mean? (\textquotedblleft{}A fort\textquotedblright{}.) Say it once more.
\end{tcolorbox}
\section[\texorpdfstring{\ml{സൈന്യം} (sainyaṁ)}{സൈന്യം (sainyaṁ)}]{\ml{സൈന്യം} (sainyaṁ) — an army}

\label{lesson:ML-C399-sainyam}



\begin{quote}
Before the new one: say the Malayalam for an empire, then the Malayalam for a fort.
\end{quote}

\subsection*{You'll want to know: \ml{സൈന്യം}}
\textbf{\ml{സൈന്യം}} — \emph{sainyaṁ} — \textquotedblleft{}an army\textquotedblright{}.

\begin{quote}
\ml{സൈന്യം} \ml{അതിർത്തിയിൽ} \ml{എത്തി}.
\end{quote}

\begin{quote}
\emph{sainyaṁ atirttiyil etti.}
\end{quote}

\begin{quote}
The army reached the border.
\end{quote}

The troops of a country; one soldier is a \ml{സൈനികൻ}.

\begin{cousinweb}
\textbf{\ml{സൈന്യം}} (\emph{sainyaṁ}) is borrowed from Sanskrit \emph{sainya}, \textquotedblleft{}belonging to an army\textquotedblright{}, and as a noun \textquotedblleft{}an army\textquotedblright{} (CDIAL 13602).
\end{cousinweb}

\begin{grammarlens}[title={What you've built}]
One more word for telling what happened.
\end{grammarlens}

\subsection*{Your turn}
\begin{itemize}
\raggedright
  \item \emph{Say it:} \emph{sainyaṁ}
  \item \emph{Say it:} \emph{sainyaṁ}, once more
  \item \emph{Say it:} \emph{sainyaṁ atirttiyil etti.} — The army reached the border.
  \item \emph{Recall:} say the Malayalam for an empire, then the Malayalam for a fort, then say \emph{sainyaṁ} again
\end{itemize}

\begin{tcolorbox}[breakable,colback=teal!4,colframe=teal!35!black,title={Before you move on}]
What does \ml{സൈന്യം} mean? (\textquotedblleft{}An army\textquotedblright{}.) Say it once more.
\end{tcolorbox}
\section[\texorpdfstring{\ml{പട്ടാളക്കാരൻ} (paṭṭāḷakkāran)}{പട്ടാളക്കാരൻ (paṭṭāḷakkāran)}]{\ml{പട്ടാളക്കാരൻ} (paṭṭāḷakkāran) — a soldier}

\label{lesson:ML-C399-pattalakkaran}



\begin{quote}
Before the new one: say the Malayalam for a fort, then the Malayalam for an army.
\end{quote}

\subsection*{You'll want to know: \ml{പട്ടാളക്കാരൻ}}
\textbf{\ml{പട്ടാളക്കാരൻ}} — \emph{paṭṭāḷakkāran} — \textquotedblleft{}a soldier\textquotedblright{}.

\begin{quote}
\ml{എന്റെ} \ml{മുത്തച്ഛൻ} \ml{ഒരു} \ml{പട്ടാളക്കാരൻ} \ml{ആയിരുന്നു}.
\end{quote}

\begin{quote}
\emph{enṟe muttacchan oru paṭṭāḷakkāran āyirunnu.}
\end{quote}

\begin{quote}
My grandfather was a soldier.
\end{quote}

Plural: \ml{പട്ടാളക്കാർ}, \emph{paṭṭāḷakkār}.

\ml{പട്ടാളം}, the army, + -\ml{ക്കാരൻ}.

\begin{cousinweb}
\textbf{\ml{പട്ടാളക്കാരൻ}} (\emph{paṭṭāḷakkāran}) is built from \ml{പട്ടാളം}, \textquotedblleft{}the troops, the army\textquotedblright{}, whose origin is not settled, and the ending -\ml{ക്കാരൻ}, \textquotedblleft{}one who belongs to it\textquotedblright{}.
\end{cousinweb}

\begin{grammarlens}[title={What you've built}]
One more word for telling what happened.
\end{grammarlens}

\subsection*{Your turn}
\begin{itemize}
\raggedright
  \item \emph{Say it:} \emph{paṭṭāḷakkāran}
  \item \emph{Say it:} \emph{paṭṭāḷakkāran}, once more
  \item \emph{Say it:} \emph{enṟe muttacchan oru paṭṭāḷakkāran āyirunnu.} — My grandfather was a soldier.
  \item \emph{Recall:} say the Malayalam for a fort, then the Malayalam for an army, then say \emph{paṭṭāḷakkāran} again
\end{itemize}

\begin{tcolorbox}[breakable,colback=teal!4,colframe=teal!35!black,title={Before you move on}]
What does \ml{പട്ടാളക്കാരൻ} mean? (\textquotedblleft{}A soldier\textquotedblright{}.) Say it once more.
\end{tcolorbox}
\section[\texorpdfstring{\ml{കീഴടക്കുക} (kīḻaṭakkuka)}{കീഴടക്കുക (kīḻaṭakkuka)}]{\ml{കീഴടക്കുക} (kīḻaṭakkuka) — to conquer}

\label{lesson:ML-C399-kilatakkuka}



\begin{quote}
Before the new one: say the Malayalam for an army, then the Malayalam for a soldier.
\end{quote}

\subsection*{You'll want to know: \ml{കീഴടക്കുക}}
\textbf{\ml{കീഴടക്കുക}} — \emph{kīḻaṭakkuka} — \textquotedblleft{}to conquer\textquotedblright{}.

\begin{quote}
\ml{രാജാവ്} \ml{അടുത്ത} \ml{രാജ്യം} \ml{കീഴടക്കി}.
\end{quote}

\begin{quote}
\emph{rājāvŭ aṭutta rājyaṁ kīḻaṭakki.}
\end{quote}

\begin{quote}
The king conquered the neighbouring country.
\end{quote}

Past: \ml{കീഴടക്കി}, \emph{kīḻaṭakki}.

\ml{കീഴ്}, below, + \ml{അടക്കുക}, to press down: to bring under.

\begin{cousinweb}
\textbf{\ml{കീഴടക്കുക}} (\emph{kīḻaṭakkuka}) is a native Dravidian compound of \ml{കീഴ്}, \textquotedblleft{}below, under\textquotedblright{} (DEDR 1619), and \ml{അടക്കുക}, \textquotedblleft{}to press down, to subdue\textquotedblright{} (DEDR 63).
\end{cousinweb}

\begin{grammarlens}[title={What you've built}]
One more word for telling what happened.
\end{grammarlens}

\subsection*{Your turn}
\begin{itemize}
\raggedright
  \item \emph{Say it:} \emph{kīḻaṭakkuka}
  \item \emph{Say it:} \emph{kīḻaṭakkuka}, once more
  \item \emph{Say it:} \emph{rājāvŭ aṭutta rājyaṁ kīḻaṭakki.} — The king conquered the neighbouring country.
  \item \emph{Recall:} say the Malayalam for an army, then the Malayalam for a soldier, then say \emph{kīḻaṭakkuka} again
\end{itemize}

\begin{tcolorbox}[breakable,colback=teal!4,colframe=teal!35!black,title={Before you move on}]
What does \ml{കീഴടക്കുക} mean? (\textquotedblleft{}To conquer\textquotedblright{}.) Say it once more.
\end{tcolorbox}
